\documentclass[10pt,a4paper]{amsart}
\usepackage[margin=19mm]{geometry}
\usepackage{amsmath,amssymb,mathtools}
\usepackage[scr=rsfs,cal=euler]{mathalfa}
\usepackage{xfrac}
\usepackage[unicode,hidelinks,bookmarks=false]{hyperref}
\newcommand{\ie}{i.e.\ }
\newcommand{\cf}{cf.\ }
\newcommand{\resp}{respectively\ }
\newcommand{\Subsection}{\subsection}
\providecommand{\nobreakdash}{\nobreak\hbox{-}}
\setlength{\emergencystretch}{2em}
\multlinegap0pt
\usepackage{amssymb}
\usepackage{xcolor}
\usepackage{enumerate}
\usepackage{tikz}
\newtheorem{lem}{Lemma}
\newcommand{\Diff}{\mathrm{Diff}}
\title{Local centralizer rigidity for a non-generic diagonal map}
\author{Zhijing Wendy Wang, Amie Wilkinson}
\date{Source: arXiv:2609.04643v1 (CC BY 4.0)}
\begin{document}
\maketitle

\noindent\emph{Source and licence.} Local centralizer rigidity for a non-generic diagonal map. Version arXiv:2609.04643v1 (CC BY 4.0), distributed by arXiv under the Creative Commons Attribution 4.0 International licence, http://creativecommons.org/licenses/by/4.0/, as recorded in the arXiv licence metadata for this exact version and shown on the abstract page https://arxiv.org/abs/2609.04643v1. Full licence notice: \texttt{licenses/arxiv-2609.04643-CC-BY-4.0.txt}.\par\smallskip

\noindent\textit{Editorial scope.} Counted unit: the article's lem lem:app-relative-ac (source0.tex lines 3497--3506). The article's complete original proof of it is delivered with it (source0.tex lines 3508--3601, one \texttt{proof} environment of 94 lines). No mathematics is written, repaired or re-derived: the statement and the proof are the article's own lines, in the article's own notation, and no retained line is substituted or re-flowed.  Retained uncounted as the source-local context the proof needs: lines 3348--3369; lines 3480--3484.  Nothing was omitted from any retained span, so the package carries no omission record.  No retained line contains a citation, so the wrapper carries no bibliography and no bibliography entry.  No retained line contains a figure environment, an image-inclusion command or drawing code, so no image is delivered and the package declares no asset dependency.  Editorial formatting only, in addition: an AMS \texttt{amsart} preamble that restates the article's own declarations and macros used by the retained lines, namely the environments \texttt{lem} on separate counters, together with the standard packages those lines need.  The retained environments print in this document as Lemma 1 in order of appearance, whereas the article prints the same items with the numbering of its own full document.  The complete native source of the article as distributed in the arXiv e-print for this exact version is bundled unchanged as \texttt{sources/209379/source0.tex}.
\begin{lem}
\label{lem:app-relative-pesin}
Let $T\in\Diff^2(M)$ preserve a uniformly $C^2$ foliation $\mathcal F$,
and let $\nu$ be a $T$-invariant probability measure.  Assume that
$DT|_{T\mathcal F}$ has no zero Lyapunov exponent at $\nu$-almost every
point.  There is a countable family of compact Pesin blocks
$\{\Lambda_j\}$ covering the regular set modulo a $\nu$-null set such that, on
each $\Lambda_j$, the following data are uniform:
\begin{enumerate}
\item the dimensions of $E^-_{\mathcal F,T}$ and $E^+_{\mathcal F,T}$;
\item local $C^1$ disks
$\mathcal W^-_{\mathcal F,T,\mathrm{loc}}(x)$ and
$\mathcal W^+_{\mathcal F,T,\mathrm{loc}}(x)$ tangent at $x$ to
$E^-_{\mathcal F,T}(x)$ and $E^+_{\mathcal F,T}(x)$;
\item a radius $r_j>0$ and constants $C_j>0$, $\gamma_j>0$ such that the
relative stable and unstable disks have intrinsic radius at least $r_j$ and
satisfy the usual exponential contraction estimates.
\end{enumerate}
In particular, these disks are contained in the relative stable and unstable
sets $\mathcal W^-_{\mathcal F,T}(x)$ and
$\mathcal W^+_{\mathcal F,T}(x)$ defined above.
\end{lem}

\begin{equation}\label{eq:leafwise-stable-contraction}
 d_{\mathcal F}(T^ny,T^nx)
 \le C_\Lambda e^{-\gamma_\Lambda n}d_{\mathcal F}(y,x),
 \qquad n\ge0,
\end{equation}

\begin{lem}
\label{lem:app-relative-ac}
On the blocks from Lemma~\ref{lem:app-relative-pesin}, the local relative
stable and unstable laminations have absolutely continuous holonomy inside
$\mathcal F$.  More precisely, fix two blocks $\Lambda,\Lambda'$ and a
foliation box.  On the part of a relative stable holonomy whose initial
point lies in $\Lambda$ and terminal point lies in $\Lambda'$, the
leafwise measures of a set and of its holonomy image are comparable, with
constants depending only on the two blocks and the foliation box.
\end{lem}

\begin{proof}
We give the stable case.  Work inside one $\mathcal F$-plaque and let
$\Sigma_0,\Sigma_1$ be sufficiently small smooth disks transverse to the
relative stable plaques.  Let
\[
 h:D_0\subset\Sigma_0\longrightarrow D_1\subset\Sigma_1
\]
be the relative stable holonomy.  Fix two blocks $\Lambda,\Lambda'$ and
restrict to
\[
 D_{\Lambda,\Lambda'}
 =\{z\in D_0:z\in\Lambda,\ h(z)\in\Lambda'\}.
\]
Since there are only countably many blocks, it is enough to obtain a uniform
absolute-continuity estimate on each such measurable piece.  After a
countable further decomposition, all charts, cone bounds, and local
extensions used below may be taken uniformly on the piece.

After shrinking the foliation box, the tangent spaces of $\Sigma_0$ and
$\Sigma_1$ lie in a common cone transverse to the stable direction.  Forward
graph transform narrows this cone.  At the same time, if $z$ and $h(z)$ are
paired by stable holonomy, then their forward iterates approach one another
at the exponential rate in \eqref{eq:leafwise-stable-contraction}.  Hence,
in the time-$n$ Lyapunov chart, the relevant pieces of
$T^n\Sigma_0$ and $T^n\Sigma_1$ are $C^1$-close graphs over the same
$E^+$-coordinate plane.  Let
\[
 \pi_n:T^n\Sigma_0\longrightarrow T^n\Sigma_1
\]
be the map obtained by matching the $E^+$ coordinates.  Then
\[
 \|D\pi_n-I\|\longrightarrow0
\]
uniformly on each fixed measurable chart piece.

Define the approximate holonomy
\[
 h_n=(T^n|_{\Sigma_1})^{-1}\circ\pi_n\circ(T^n|_{\Sigma_0}).
\]
Then $h_n\to h$ uniformly.  For $z\in D_{\Lambda,\Lambda'}$, put
\[
 w_n=h_n(z),\qquad z_k=T^kz,\qquad w_{k,n}=T^kw_n,
\]
and
\[
 P_{0,k}=DT^k(T_z\Sigma_0),
 \qquad
 P_{1,k,n}=DT^k(T_{w_n}\Sigma_1).
\]
The stable contraction, the forward and backward cone estimates, and the
$C^1$ regularity of $DT|_{T\mathcal F}$ give constants $C,c>0$, depending
only on the two blocks and the chosen chart piece, such that
\begin{equation}\label{eq:leafwise-distortion}
 \left|
 \log\bigl|\det(DT(z_k)|_{P_{0,k}})\bigr|
 -
 \log\bigl|\det(DT(w_{k,n})|_{P_{1,k,n}})\bigr|
 \right|
 \le C\bigl(e^{-ck}+e^{-c(n-k)}\bigr)
\end{equation}
for $0\le k<n$.  The first term compares the two true stable orbits, while
the second accounts for pulling backward the short projection made at time
$n$.

Summing \eqref{eq:leafwise-distortion} and using the uniform Jacobian bounds
for $\pi_n$ and $\pi_n^{-1}$ gives
\begin{equation}\label{eq:leafwise-holonomy-jacobian}
 C_1^{-1}\le\operatorname{Jac}h_n\le C_1,
\end{equation}
with the same estimate for $h_n^{-1}$ and with $C_1$ independent of $n$.
By the area formula, for every Borel set
$B\subset D_{\Lambda,\Lambda'}$,
\[
 \operatorname{vol}_{\Sigma_1}(h_n(B))
 \le C_1\operatorname{vol}_{\Sigma_0}(B).
\]
Interchanging $\Sigma_0$ and $\Sigma_1$ gives the corresponding
estimate for $h_n^{-1}$ and also $h_n^{-1}\to h^{-1}$ uniformly.  Thus
\[
 (h_n)_*\operatorname{vol}_{\Sigma_0}
 \le C_1\operatorname{vol}_{\Sigma_1}
\]
on the measurable piece, and similarly for the inverse push-forward.
Passing to weak limits and using the uniform convergence $h_n\to h$ gives
\[
 \operatorname{vol}_{\Sigma_1}(h(B))
 \le C_1\operatorname{vol}_{\Sigma_0}(B).
\]
Thus $h$ sends null sets to null sets on
$D_{\Lambda,\Lambda'}$.  Taking the countable union over pairs of blocks,
chart pieces, transversals, and foliation boxes proves absolute continuity
of the relative stable lamination inside $\mathcal F$.  The unstable case
follows by applying the same argument to $T^{-1}$.
\end{proof}

\end{document}
